\documentclass[10pt,a4paper]{amsart}
\usepackage[margin=19mm]{geometry}
\usepackage{amsmath,amssymb,mathtools}
\usepackage[scr=rsfs,cal=euler]{mathalfa}
\usepackage{xfrac}
\usepackage[unicode,hidelinks,bookmarks=false]{hyperref}
\newcommand{\ie}{i.e.\ }
\newcommand{\cf}{cf.\ }
\newcommand{\resp}{respectively\ }
\newcommand{\Subsection}{\subsection}
\providecommand{\nobreakdash}{\nobreak\hbox{-}}
\setlength{\emergencystretch}{2em}
\multlinegap0pt
\usepackage{amssymb}
\usepackage{xcolor}
\usepackage[shortlabels]{enumitem}
\usepackage{float}
\usepackage{tikz}
\newtheorem{lemma}{Lemma}
\title{Crossed products as compact quantum metric spaces}
\author{Mario Klisse}
\date{Source: arXiv:2303.17903v2 (CC BY 4.0)}
\begin{document}
\maketitle

\noindent\emph{Source and licence.} Crossed products as compact quantum metric spaces. Version arXiv:2303.17903v2 (CC BY 4.0), distributed by arXiv under the Creative Commons Attribution 4.0 International licence, http://creativecommons.org/licenses/by/4.0/, as recorded in the arXiv licence metadata for this exact version and shown on the abstract page https://arxiv.org/abs/2303.17903v2. Full licence notice: \texttt{licenses/arxiv-2303.17903-CC-BY-4.0.txt}.\par\smallskip

\noindent\textit{Editorial scope.} Counted unit: the article's lemma ConditionalExpectation-1 (source0.tex lines 757--763). The article's complete original proof of it is delivered with it (source0.tex lines 765--816, one \texttt{proof} environment of 52 lines). No mathematics is written, repaired or re-derived: the statement and the proof are the article's own lines, in the article's own notation, and no retained line is substituted or re-flowed.  Retained uncounted as the source-local context the proof needs: lines 614--624.  Nothing was omitted from any retained span, so the package carries no omission record.  No retained line contains a citation, so the wrapper carries no bibliography and no bibliography entry.  No retained line contains a figure environment, an image-inclusion command or drawing code, so no image is delivered and the package declares no asset dependency.  Editorial formatting only, in addition: an AMS \texttt{amsart} preamble that restates the article's own declarations and macros used by the retained lines, namely the environments \texttt{lemma} on separate counters, together with the standard packages those lines need.  The retained environments print in this document as Lemma 1 in order of appearance, whereas the article prints the same items with the numbering of its own full document.  The complete native source of the article as distributed in the arXiv e-print for this exact version is bundled unchanged as \texttt{sources/139489/source0.tex}.
\begin{lemma} \label{Commutator identity} Let $\alpha:G\rightarrow\text{Aut}(A)$
be an action of a discrete group $G$ on a separable unital C$^{\ast}$-algebra
$A\subseteq\mathcal{B}(\mathcal{H}_{A})$ and let $\ell:G\rightarrow\mathbb{R}_{+}$
be a proper length function on $G$. For every $x=\sum_{g\in G}a_{g}\lambda_{g}\in C_{c}(G,A)\subseteq\mathcal{B}(\mathcal{H})$
with $(a_{g})_{g\in G}\subseteq A$, $\mathcal{H}:=\mathcal{H}_{A}\otimes\ell^{2}(G)$,
\[
[1\otimes M_{\ell},x]=\sum_{g\in G}(1\otimes\varphi_{g}^{\ell})a_{g}\lambda_{g}
\]
where $\varphi_{g}^{\ell}$ is viewed as a multiplication operator
$\delta_{h}\mapsto\varphi_{g}^{\ell}(h)$ in $\ell^{\infty}(G)\subseteq\mathcal{B}(\ell^{2}(G))$.
\end{lemma}

\begin{lemma} \label{ConditionalExpectation-1} Let $H\subseteq G$
be a subgroup. Then there exists a contractive linear map $\mathbb{E}_{H}$
on $\mathcal{B}(\mathcal{H})$ such that for every $x=\sum_{g\in G}a_{g}\lambda_{g}\in C_{c}(G,\mathcal{A})\subseteq\mathcal{B}(\mathcal{H})$
with $(a_{g})_{g\in G}\subseteq\mathcal{A}$, the identities
$\mathbb{E}_{H}(x)=\sum_{h\in H}a_{h}\lambda_{h} \in C_c(H,\mathcal{A})$, $\mathbb{E}_{H}([D_{A}\otimes1,x]\lambda_{g^{-1}})\lambda_{g}=[D_{A}\otimes1,\mathbb{E}_{H}(x\lambda_{g^{-1}})\lambda_{g}]$
and $\mathbb{E}_{H}([1\otimes M_{\ell},x]\lambda_{g^{-1}})\lambda_{g}=[1\otimes M_{\ell},\mathbb{E}_{H}(x\lambda_{g^{-1}})\lambda_{g}]$
hold. \end{lemma}

\begin{proof} Let $(g_{i})_{i\in I}\subseteq G$ be a family of elements
with $G=\bigcup_{i\in I}Hg_{i}$ and $Hg_{i}\neq Hg_{j}$ for $i\neq j$.
For $i\in I$, write $P_{Hg_{i}}$ for the orthogonal projection onto
the closed subspace of $\ell^{2}(G)$ spanned by all orthonormal basis
vectors $\delta_{hg_{i}}$, $h\in H$. We claim that the linear map
$\mathbb{E}_{H}$ given by $x\mapsto\sum_{i\in I}(1\otimes P_{Hg_{i}})x(1\otimes P_{Hg_{i}})$
satisfies the required conditions, where the sum converges in the strong
operator topology. Indeed, for every $x\in\mathcal{B}(\mathcal{H})$,
\[
\Vert\sum_{i\in I}(1\otimes P_{Hg_{i}})x(1\otimes P_{Hg_{i}})\Vert\leq\sup_{i\in I}\left\Vert (1\otimes P_{Hg_{i}})x(1\otimes P_{Hg_{i}})\right\Vert \leq\left\Vert x\right\Vert ,
\]
as the operators $(1\otimes P_{Hg_{i}})x(1\otimes P_{Hg_{i}})$, $i\in I$
have pairwise orthogonal support and ranges. For $x=\sum_{g\in G}a_{g}\lambda_{g}\in C_{c}(G,\mathcal{A})\subseteq\mathcal{B}(\mathcal{H})$
with $(a_{g})_{g\in G}\subseteq\mathcal{A}$ and $\xi\in\mathcal{H}_{A}$,
$h^{\prime}\in H$, $i\in I$, we further find 
\begin{eqnarray}
\nonumber
\left(\mathbb{E}_{H}(x)\right)(\xi\otimes\delta_{h^{\prime}g_{i}}) &=& ((1\otimes P_{Hg_{i}})\sum_{g\in G}a_{g}\lambda_{g})(\xi\otimes\delta_{h^{\prime}g_{i}}) \\
\nonumber
&=& \sum_{g\in G}(1\otimes P_{Hg_{i}})\left(((gh^{\prime}g_{i})^{-1}.a_{g})\xi\otimes\delta_{gh^{\prime}g_{i}}\right) \\
\nonumber
&=& \sum_{h\in H}\left(((hh^{\prime}g_{i})^{-1}.a_{h})\xi\otimes\delta_{hh^{\prime}g_{i}}\right) \\
\nonumber
&=& (\sum_{h\in H}a_{h}\lambda_{h})(\xi\otimes\delta_{h^{\prime}g_{i}}),
\end{eqnarray}
so that $\mathbb{E}_{H}(x)=\sum_{h\in H}a_{h}\lambda_{h}$. For $g,g^{\prime}\in G$,
there exist $i\in I$ and $h^{\prime}\in H$ such that $gg^{\prime}=h^{\prime}g_{i}$.
It follows that
\begin{eqnarray}
\nonumber
\left(\mathbb{E}_{H}([D_{A}\otimes1,x]\lambda_{g^{-1}})\lambda_{g}\right)(\xi\otimes\delta_{g^{\prime}}) &=& ((1\otimes P_{Hg_{i}})\sum_{g^{\prime\prime}\in G}[D_{A}\otimes1,a_{g^{\prime\prime}}]\lambda_{g^{\prime\prime}g^{-1}})(\xi\otimes\delta_{h^{\prime}g_{i}}) \\
\nonumber
&=& (\sum_{g^{\prime\prime}\in Hg}[D_{A}\otimes1,a_{g^{\prime\prime}}]\lambda_{g^{\prime\prime}g^{-1}})(\xi\otimes\delta_{h^{\prime}g_{i}}) \\
\nonumber
&=& (\sum_{g^{\prime\prime}\in Hg}[D_{A}\otimes1,a_{g^{\prime\prime}}]\lambda_{g^{\prime\prime}})(\xi\otimes\delta_{g^{\prime}}) \\
\nonumber
&=& \left([D_{A}\otimes1,\mathbb{E}_{H}(x\lambda_{g^{-1}})\lambda_{g}]\right)(\xi\otimes\delta_{g^{\prime}})
\end{eqnarray}
and with Lemma \ref{Commutator identity},
\begin{eqnarray}
\nonumber
\left(\mathbb{E}_{H}([1\otimes M_{\ell},x]\lambda_{g^{-1}})\lambda_{g}\right)(\xi\otimes\delta_{g^{\prime}}) &=& ((1\otimes P_{Hg_{i}})\sum_{g\in G}(1\otimes\varphi_{g^{\prime\prime}}^{\ell})a_{g^{\prime\prime}}\lambda_{g^{\prime\prime}g^{-1}})(\xi\otimes\delta_{h^{\prime}g_{i}}) \\
\nonumber
&=& (\sum_{g^{\prime\prime}\in Hg}(1\otimes\varphi_{g^{\prime\prime}}^{\ell})a_{g^{\prime\prime}}\lambda_{g^{\prime\prime}g^{-1}})(\xi\otimes\delta_{h^{\prime}g_{i}}) \\
\nonumber
&=& (\sum_{g^{\prime\prime}\in Hg}(1\otimes\varphi_{g^{\prime\prime}}^{\ell})a_{g^{\prime\prime}}\lambda_{g^{\prime\prime}})(\xi\otimes\delta_{g^{\prime}}) \\
\nonumber
&=& ([1\otimes M_{\ell},\mathbb{E}_{H}(x\lambda_{g^{-1}})\lambda_{g}])(\xi\otimes\delta_{g^{\prime}}).
\end{eqnarray}
We deduce that $\mathbb{E}_{H}([D_{A}\otimes1,x]\lambda_{g^{-1}})\lambda_{g}=[D_{A}\otimes1,\mathbb{E}_{H}(x\lambda_{g^{-1}})\lambda_{g}]$
and $\mathbb{E}_{H}([1\otimes M_{\ell},x]\lambda_{g^{-1}})\lambda_{g}=[1\otimes M_{\ell},\mathbb{E}_{H}(x\lambda_{g^{-1}})\lambda_{g}]$,
as claimed. \end{proof}

\end{document}
